\clearpage\fancyhead[L]{\small COMET.}\hypertarget{native-378-lower}{}\label{native-378-lower}
\begin{center}\textbf{COMET.}\par\vspace{4pt}\hrule\vspace{4pt}\fontsize{10}{12}\selectfont Göttingische gelehrte Anzeigen. No. 19. pp. 185\;..\;189. 31 January 1814.\par\vspace{4pt}\hrule\vspace{4pt}\end{center}

\par\indent From a letter of 30 December last year from Professor \textsc{Bessel} in Königsberg to Professor \textsc{Gauss}, we give a few details that the astronomical readers of our journal will be glad to learn promptly. Various reports have already shown that there was some ground for the prospect, mentioned several times in this journal too (see 1811, no. 130, and 1812, no. 67 [pp. 327 and 349 of this volume]), that the great comet of 1811 might be seen and observed again during the summer with excellent instruments and under favorable conditions. But the astronomers at the European observatories, despite their careful efforts, did not have that good fortune. This was partly because the twilight at night is too bright at that time of year and partly because of the less favorable climate that most of them have to contend with. From the Caucasus, from \emph{Neu-Tscherkask}, astronomers have received unique observations of this remarkable comet. Mr. von \textsc{Wisniwsky}, who has been traveling in those remote regions for several years, made them nearly a year after the comet’s passage through perihelion. Professor \textsc{Bessel}, whose contributions to the theory of this comet in more than one respect are well known and who had calculated an elliptic orbit for it, \SourceLeaf{379}received these important observations from Petersburg. He will use them later to refine the orbit as closely as possible. The following preliminary results, however, deserve to be made known now.

\Needspace{14\baselineskip}\par\indent Mr. von \textsc{Wisniwsky} observed the comet from 8 to 27 August. In that time he obtained 29 transits through the ring micrometer, whose field he determined with great precision using several pairs of stars. He compared the comet with the following stars:

\begin{center}\fontsize{10}{12}\selectfont\setlength{\tabcolsep}{5pt}\renewcommand{\arraystretch}{1.15}\begin{tabular}{c|rr|rr}
& \multicolumn{2}{c|}{Right ascension} & \multicolumn{2}{c}{Southern declination}\\\hline
I & $324^\circ$ & $34'$ & $24^\circ$ & $31'$\\
II & 323 & 10 & 24 & 59\\
III & 322 & 37 & 25 & 16\\
IV & 322 & 45 & 25 & 57\\
V & 322 & 17 & 26 & 3\\
\end{tabular}\end{center}

\Needspace{15\baselineskip}\par\indent Professor \textsc{Bessel}’s reduction gives the following positions of the comet relative to these stars:

\begin{center}\fontsize{9}{11}\selectfont\setlength{\tabcolsep}{3pt}\renewcommand{\arraystretch}{1.15}\begin{tabular}{l|c|c|c|rr|rr}
1812 & \shortstack{Mean time in\\Neu-Tscherkask} & \shortstack{Comparison\\star} & \shortstack{No. of\\obs.} & \multicolumn{2}{c|}{\shortstack{Right\\ascension}} & \multicolumn{2}{c}{Declination}\\\hline
Aug. 8 & $12^{\mathrm u}\;18^{\mathrm m}\;51^{\mathrm s}4$ & I & 4 & $-5'$ & $30''5$ & $-1'$ & $19''3$\\
11 & $12\;45\;21.7$ & II & 2 & $+25$ & 23.7 & $-3$ & 52.2\\
12 & $11\;42\;25.3$ & II & 7 & $+9$ & 7.5 & $-51$ & 15.3\\
12 & $12\;41\;41.0$ & III & 3 & $+41$ & 10.0 & $+1$ & 15.8\\
15 & $12\;10\;7.9$ & IV & 3 & $-18$ & 9.0 & $+11$ & 50.0\\
15 & $12\;30\;18.7$ & V & 5 & $+8$ & 39.8 & $+13$ & 0.0\\
17 & $13\;34\;8.4$ & V & 6 & $-25$ & 36.7 & $-1$ & 1.7\\
\end{tabular}\end{center}

\Needspace{12\baselineskip}\par\indent Only three of the stars themselves (II, IV, V) appear in the \emph{Histoire céleste}. Professor \textsc{Bessel} reduced the recorded positions to the year 1812 using his own reduction elements and found:

\begin{center}\fontsize{10}{12}\selectfont\setlength{\tabcolsep}{4pt}\renewcommand{\arraystretch}{1.15}\begin{tabular}{c|rrr|rrr}
& \multicolumn{3}{c|}{Right ascension} & \multicolumn{3}{c}{Southern declination}\\\hline
II & $323^\circ$ & $9'$ & $1''5$ & $24^\circ$ & $59'$ & $26''5$\\
IV & 322 & 44 & 32.3 & 25 & 56 & 53.4\\
V & 322 & 16 & 17.1 & 26 & 3 & 36.9\\
\end{tabular}\end{center}

\Needspace{7\baselineskip}\par\indent He also determined the right ascensions of II and IV from observations of his own:

\begin{center}\fontsize{10}{12}\selectfont\setlength{\tabcolsep}{4pt}\renewcommand{\arraystretch}{1.15}\begin{tabular}{lcrl}
II & $.\;.\;.\;.\;.\;.\;.\;.$ & $323^\circ\;9'\;5''3$ & (1 obs.)\\
IV & $.\;.\;.\;.\;.\;.\;.\;.$ & $322\;44\;56.1$ & (2 obs.)\\
\end{tabular}\end{center}

\Needspace{13\baselineskip}\par\indent Taking these determinations as the preferred basis gave the following observed positions of the comet:

\begin{center}\fontsize{9}{11}\selectfont\setlength{\tabcolsep}{3pt}\renewcommand{\arraystretch}{1.15}\begin{tabular}{l|c|rrr|rrr}
& Mean time & \multicolumn{3}{c|}{Right ascension} & \multicolumn{3}{c}{Southern declination}\\\hline
Aug. 11 & $12^{\mathrm u}\;45^{\mathrm m}\;21^{\mathrm s}7$ & $323^\circ$ & $35'$ & $9''9$: & $25^\circ$ & $3'$ & $0''4$\\
12 & $11\;42\;25.3$ & 323 & 18 & 53.7 & 25 & 14 & 59.5\\
15 & $12\;10\;7.9$ & 322 & 27 & 7.9 & 25 & 44 & 45.3\\
15 & $12\;30\;18.7$ & 322 & 25 & 37.9 & 25 & 45 & 19.0\\
17 & $13\;34\;8.4$ & 321 & 51 & 21.4 & 26 & 4 & 20.7\\
\end{tabular}\end{center}

\Needspace{13\baselineskip}\par\indent Comparison of these observations with the purely elliptic elements gave Professor \textsc{Bessel} the following differences
\SourceLeaf{380}\begin{center}\fontsize{10}{12}\selectfont\setlength{\tabcolsep}{8pt}\renewcommand{\arraystretch}{1.15}\begin{tabular}{l|c|c}
& \shortstack{Right\\ascension} & Declination\\\hline
Aug. 11 & $-31''2$ & $-10''0$:\\
12 & $-55.9$ & $+113.9$\\
15 & $-86.4$ & $+76.1$\\
15 & $-10.9$ & $+101.6$\\
17 & $-51.2$ & $+57.8$\\
\end{tabular}\end{center}

\par\indent This agreement, after an interval of more than three quarters of a year, proves how well Mr. \textsc{Bessel} knew how to use the observations. In particular, those he made himself with an excellent heliometer were exceptionally accurate. Indeed, Mr. \textsc{Bessel} would have obtained somewhat better agreement had he thought it worth taking the trouble at the time to refine his elements once more using the latest observations then available, since they had already begun to depart somewhat from the elements in the same direction. It is also clear that the perturbations can have had only a very small effect on the comet’s motion. Mr. \textsc{Bessel} is now working on the calculations for this comet.
